% K1--K4 eredményösszefoglaló, első táblázatos változat.
\setlength{\tabcolsep}{4pt}
\renewcommand{\arraystretch}{1.08}
\begin{tabularx}{\textwidth}{@{}
>{\hsize=.64\hsize\linewidth=\hsize\raggedright\arraybackslash}X
>{\hsize=1.40\hsize\linewidth=\hsize\raggedright\arraybackslash}X
>{\hsize=.80\hsize\linewidth=\hsize\raggedright\arraybackslash}X
>{\hsize=1.16\hsize\linewidth=\hsize\raggedright\arraybackslash}X@{}}
\toprule
\textbf{Kérdés} & \textbf{Rövid válasz} & \textbf{Bizonyíték} &
\textbf{Erősség és korlát} \\
\midrule
\textbf{K1.} Elkészítési idő &
40,15 $\to$ 18,86 perc: 53,0\%-os csökkenés; mind a 10 alkalmazásban
alacsonyabb átlag. Munkamenetenként átlagolva 56,0\%. &
\ref{sec:fo-eredmeny}--\ref{sec:alkalmazas-elteresek}.~szakasz;
\ref{tab:alkalmazasok}.~táblázat. &
Következetes, leíró eredmény; egyenetlen, néhol alacsony ismétlésszám. \\
\addlinespace[.5em]
\textbf{K2.} Hibák és sikerjelzések &
39 AGATA-futtatásban 353 parancshiba; a teljes mintában
103/110 teljes ötsoros sikerjelzés, öt ismert GUI-megkerüléssel. &
\ref{sec:siker-ertelem}.~szakasz;
\ref{tab:parancshibak}.~táblázat. &
Célzott audit; a sikerjelzés önmagában nem igazolja a GUI-útvonal betartását. \\
\addlinespace[.5em]
\textbf{K3.} Újabb modellek &
AGATA-val mindháromnál 3/3 teljes sikerjelzés.
GPT-6.1 Sol: 6,80 perc, 66,5\%-kal kevesebb a saját referenciánál. &
Közös Office-minta;
\ref{sec:modell-eredmeny}.~szakasz;
\ref{tab:modellek}.~táblázat. &
Előzetes: AGATA-val alkalmazásonként $n=1$;
a modellváltozat és a gondolkodási szint is eltért. \\
\addlinespace[.5em]
\textbf{K4.} MI-eszköz\-készlet &
MI-eszközkészlet: 24,92 perc; AGATA: 12,19 perc.
Az MI-eszközkészletnél 9/9 teljes sikerjelzés, három GUI-megkerüléssel. &
\ref{sec:office-eredmeny}.~szakasz;
\ref{tab:office-mi-tartomany}.~táblázat. &
Előzetes: három Office-alkalmazás; az MI-eszközkészletnél
alkalmazásonként $n=3$. \\
\bottomrule
\end{tabularx}
